\documentclass[12pt]{article}

% --- Языки и шрифты для XeLaTeX/LuaLaTeX ---
\usepackage{fontspec}
\usepackage{polyglossia}
\setdefaultlanguage{russian}
\setotherlanguage{english}

% Добавляет смайлики :)
\usepackage{wasysym}

% Подходящие шрифты с кириллицей (CMU есть почти везде)
\setmainfont{CMU Serif}
\setsansfont{CMU Sans Serif}
\setmonofont{CMU Typewriter Text}

% Математика и вёрстка
\usepackage{amsmath,amssymb}
\usepackage{geometry}
\geometry{a4paper,margin=2.5cm}

% Ссылки
\usepackage{hyperref}
\hypersetup{
  colorlinks=true,
  linkcolor=blue,
  urlcolor=blue,
  citecolor=blue
}

\begin{document}

\begin{center}
  {\LARGE Домашнее задание \#1}\\[4pt]
  {\large Курс: Введение в маневрирование КА}\\[2pt]
  {\small Дедлайн: \underline{06.10.2025, 23:59} }
\end{center}

\bigskip

\section*{Инструкция}
Для выполнения задания необходимо скачать на свой компьютер \href{https://github.com/Chelovechek-move/maneuvering}{этот репозиторий}. В файле \texttt{README} вы найдёте короткие инструкции для начала работы с проектом. Запустите тесты с помощью команды \texttt{pytest}. После её выполнения вы должны увидеть одну ошибку: 
\begin{center}
    \texttt{ERROR tests/maneuvers/quasi\_circular/transition/test\_coplanar.py}    
\end{center}
Если это так, значит всё в порядке \smiley. Её мы исправим в ходе выполнения заданий. 

Перед началом работы посмотрите на структуру репозитория, почитайте комментарии в файлах, потратьте на это $10$ минут. Исходный код находится в директории \texttt{src}, тесты в \texttt{tests}, исследовательские примеры, функции для построения графиков и другой код в \texttt{examples}. 

Во всех заданиях не требуется точных ответов, не нужно ничего считать. Главное объяснить смысл, можно своими словами, можно неправильно. Просто напишите свои мысли, больше ничего не надо. 

% Пример ссылки: \href{https://example.com}{кликабельный текст} или просто URL: \url{https://example.com}.

\section*{Задания}

\begin{enumerate}
  \item \textbf{Оптимальные алгоритмы компланарного маневрирования.} \\
  Найдите файлик \texttt{coplanar.py} вот тут:
  \begin{center}
      \texttt{src/maneuvering/maneuvers/quasi\_circular/transition/coplanar.py}
  \end{center}
  В нём реализованы оптимальные алгоритмы компланарного маневрирования, которые мы вывели на предыдущих лекциях. Основная структура функций уже написана. Ваша задача - найти все строки с таким текстом:
  \begin{center}
      \texttt{\# СЮДА НАДО НАПИСАТЬ КОД ...}
  \end{center}
  и написать на их месте нужные формулы. Все эти формулы могут быть найдены в презентации с номером $4$. После этого ещё раз запустите тесты командой \texttt{pytest}. Теперь всё должно заработать. Поздравляю, мы только что написали первые рабочие алгоритмы. Перейдём к их анализу.

  \item \textbf{Точность расчёта маневрирования.} \\
  Откройте файл \texttt{linear\_coplanar.ipynb} в папке \texttt{examples}. Активируйте несколько первых ячеек с импортом функций и инициализацией констант. Далее активируйте ячейки, относящимся к текущему заданию. Вы увидите $4$ графика, на каждом из которых построены две поверхности ошибок (расстояний между орбитами). Одна из них показывает расстояние между начальной и целевой орбитой $d_i = d(\textbf{Э}_i, \textbf{Э}_t)$, а другая - между финальной (полученной после выполнения маневрирования) и целевой $d_f = d(\textbf{Э}_f, \textbf{Э}_t)$. По двум другим осям отложены изменения эксцентриситета и большой полуоси. Каждый из графиков соответствует своему значению изменения аргумента перицентра. Ответьте на следующие вопросы:
  \begin{itemize}
      \item Наши алгоритмы вообще работают? Как соотносятся $d_i$ и $d_f$? 
      \item Уберите со всех графиков поверхность $d_i$, оставьте только поверхность с ошибками маневрирования $d_f$. Как влияют на значение $d_f$ параметры $\Delta e$, $\Delta a$ и $\Delta \omega$? Какой из них влияет сильнее? С чем это связано?  
      \item Почему поверхность $d_f$ не совпадает с нулём, мы же вроде решили задачу? Или нет?
  \end{itemize}
  Для ответа на вопросы можете построить и свои графики, если так будет удобнее. \textit{Подсказка: вспомните в каких приближениях была получена система терминальных условий.}

  \item \textbf{Влияние импульсов на элементы орбиты.} \\
  Продолжаем работать с файлом \texttt{linear\_coplanar.ipynb}. Запустите ячейки относящиеся к заданию $3$. Вы увидите графики зависимости кеплеровых элементов орбиты от абсолютного угла, который пролетел аппарат. Моменты скачков - прикладываемые манёвры. Ответьте на следующие вопросы:
  \begin{itemize}
      \item Как одиночный импульс влияет на изменение элементов орбиты? \textit{Подсказка: обратитесь к материалам лекции $2$.}
      
      \item Как зависят от типа оптимального маневрирования (пересекающиеся или не пересекающиеся орбиты) направления прикладываемых импульсов, а следовательно и изменения $a$ и $e$? \textit{Подсказка: вспомните графики для годографа базис-вектора.}

      \item Изменяются ли вне плоскостные параметры: $i$ и $\Omega$? Почему?

      \item Изменяется ли истинная аномалия при приложении импульсов скорости? Почему? Можем ли мы в рамках используемых алгоритмов получить некоторую целевую истинную аномалию? Можем ли мы ограничить изменения истинной аномалии? Почему истинная аномалия изменяется даже в промежутках между манёврами?
  \end{itemize}

  \item \textbf{Одноимпульсное маневрирование.} \\
  При выводе алгоритмов оптимального двухимпульсного маневрирования мы рассмотрели случай пересекающихся орбит. Как вам известно из лекции $2$, переход между двумя орбитами, имеющими общую точку, можно осуществить с использованием лишь одного импульса. \\ 
  Найдите такой одноимпульсный манёвр и точку его приложения. Рассчитайте характеристическую скорость и условия применимости такого маневрирования. Что оптимальнее: один импульс или два? Совпадают ли полученные вами условия с теми, что мы получили на лекции? \\
  \textit{Подсказка: Рассмотрите систему условий выхода на целевую орбиту для компланарного перехода. Если импульс один, то это система из трёх уравнений с тремя неизвестными. Решите эту систему.}

  \stepcounter{enumi}
  \item[\textbf{\arabic{enumi}$^\ast$}.] \textbf{Точное решение задачи перехода.} (необязательное) \\
  Как вы могли узнать при выполнении задания $2$, наши алгоритмы оптимального маневрирования решают поставленную задачу перехода с некоторой ошибкой. Иногда эта ошибка может принимать достаточно большое значение и тогда возникает вопрос, что вообще значит "решить задачу" и есть ли смысл говорить об оптимальности такого решения вообще? \\
  Давайте попробуем решить задачу перехода точно. Для этого воспользуемся несколькими фактами, которые мы уже знаем: 
  \begin{itemize}
      \item Задача перехода не фиксирует время, за которое этот переход будет выполнен (аналогом времени может служить и любая быстро меняющаяся переменная, например, истинная аномалия).
      \item Наши алгоритмы перехода всё-таки работают: ошибка между финальной и целевой орбитами меньше ошибки между начальной и целевой орбитами.
      \item Чем ближе начальная и целевая орбиты (меньше начальная ошибка между ними), тем меньшую ошибку мы в итоге получим.
  \end{itemize}
  Вследствие вышеперечисленного напрашивается следующий алгоритм:
  \begin{itemize}
      \item Рассчитать маневрирование;
      \item Выполнить манёвры;
      \item Получить некоторую промежуточную орбиту;
      \item Проверить: наша текущая орбита совпадает с целевой с достаточной точностью? Если да, то закончить расчёт. Если нет, то вернуться к первому пункту.
  \end{itemize}
  Теоретически такой алгоритм может позволить решить задачу перехода с любой наперёд заданной точностью. \\ 
  Ваша задача: реализовать описанный алгоритм. Проанализируйте, как меняется в зависимости от заданной точности время исполнения маневрирования? С какой максимальной точностью можно осуществить маневрирование? Что вы можете сказать об оптимальности рассчитанных манёвров? Постройте графики изменения элементов, как в задании $3$. Постройте графики изменения характеристической скорости. Можете ли вы предложить другие алгоритмы точного решения задачи перехода?

  
  Для решения этой задачи вам пригодятся функции:
  \begin{itemize}
      \item \texttt{distance\_orbit()} - рассчитывает расстояние между орбитами;
      \item \texttt{coplanar\_analytical()} - рассчитывает оптимальное двухимпульсное маневрирование;
      \item \texttt{execute()} - выполняет активное маневрирование, возвращает финальное состояние;
      \item \texttt{execute\_batch()} - выполняет активное маневрирование, возвращает последовательность состояний и соответствующих им абсолютных углов. Пригодится при построении графиков. Тут главное не запутаться в том, откуда отсчитываются углы в каждом из манёвров.
  \end{itemize}
  
\end{enumerate}

\end{document}
